\documentclass[11pt,a4paper]{article}
% =====================================================================
%  Shared preamble for the Part V lecture notes (lectures 19-22).
%
%  These four lectures sit outside Geron, so the notes ARE the primary
%  source and are examined on the same terms as the textbook parts.
%  Everything here is chosen to make that readable on paper: one column,
%  generous measure, and the same six-colour palette as the slides so a
%  student moving between the two is not re-learning what red means.
%
%  Build with pdflatex (twice, for the toc and the refs):
%      cd notes && latexmk -pdf lecture-19.tex
%  or  make -C notes
% =====================================================================
\usepackage[T1]{fontenc}
\usepackage[utf8]{inputenc}
\usepackage{newtxtext,newtxmath}      % Times-like: prints well, reads well
\usepackage[scaled=0.86]{inconsolata} % code, at a size that sits beside Times
\usepackage{microtype}
% newtx defines \Bbbk, \openbox, \proof and \endproof; amssymb and amsthm
% then redefine them and abort. Freeing the four names before those two
% packages load is the standard fix, and it costs nothing here: we use
% amsthm only for \newtheorem.
\let\Bbbk\relax
\usepackage{amsmath,amssymb}
\let\openbox\relax
\let\proof\relax
\let\endproof\relax
\usepackage{amsthm}
\usepackage{booktabs}
\usepackage{enumitem}
\usepackage{xcolor}
\usepackage{tcolorbox}
\usepackage{titlesec}
\usepackage{graphicx}
\usepackage[hidelinks]{hyperref}
\tcbuselibrary{skins,breakable}

% ---- the course palette, from assets/css/custom.css -------------------
\definecolor{aimlprimary}{HTML}{0B3D62}   % deep blue  - structure
\definecolor{aimlaccent} {HTML}{C0392B}   % brick red  - failures
\definecolor{aimlsuccess}{HTML}{14663A}   % green      - repairs
\definecolor{aimlmath}   {HTML}{6C3483}   % purple     - the derivation
\definecolor{aimlmuted}  {HTML}{4B5563}
\definecolor{aimlrule}   {HTML}{B0BCC7}
\definecolor{aimlsoft}   {HTML}{F4F7F9}

\usepackage[a4paper,margin=32mm,top=30mm,bottom=30mm]{geometry}
\setlength{\parskip}{0.45\baselineskip}
\setlength{\parindent}{0pt}
\linespread{1.06}

% ---- headings --------------------------------------------------------
\titleformat{\section}{\Large\bfseries\color{aimlprimary}}{\thesection}{0.7em}{}
\titleformat{\subsection}{\large\bfseries\color{aimlprimary}}{\thesubsection}{0.7em}{}
\titleformat{\subsubsection}{\bfseries\color{aimlmuted}}{}{0em}{}
\titlespacing*{\section}{0pt}{1.6\baselineskip}{0.5\baselineskip}
\titlespacing*{\subsection}{0pt}{1.2\baselineskip}{0.35\baselineskip}

% ---- boxes -----------------------------------------------------------
% One box per role, and the role is the colour. A student who has seen the
% slides already knows what each one means before reading a word of it.
\newtcolorbox{keybox}[1][]{
  breakable, enhanced, colback=aimlsoft, colframe=aimlprimary,
  boxrule=0.4pt, left=8pt, right=8pt, top=6pt, bottom=6pt,
  fonttitle=\bfseries\small, coltitle=white, colbacktitle=aimlprimary,
  attach boxed title to top left={yshift=-2mm, xshift=4mm},
  boxed title style={boxrule=0pt, arc=1pt}, #1}

\newtcolorbox{failbox}[1][]{
  breakable, enhanced, colback=white, colframe=aimlaccent,
  boxrule=0.9pt, left=8pt, right=8pt, top=6pt, bottom=6pt,
  fonttitle=\bfseries\small, coltitle=white, colbacktitle=aimlaccent,
  attach boxed title to top left={yshift=-2mm, xshift=4mm},
  boxed title style={boxrule=0pt, arc=1pt}, #1}

\newtcolorbox{derivbox}[1][]{
  breakable, enhanced, colback=white, colframe=aimlmath,
  boxrule=0.6pt, left=8pt, right=8pt, top=6pt, bottom=6pt,
  fonttitle=\bfseries\small, coltitle=white, colbacktitle=aimlmath,
  attach boxed title to top left={yshift=-2mm, xshift=4mm},
  boxed title style={boxrule=0pt, arc=1pt}, #1}

% An exercise. Deliberately the same shape as the notebook's `try` field:
% one change, and what should happen to the output.
\newtcolorbox{trybox}[1][]{
  breakable, enhanced, colback=white, colframe=aimlsuccess,
  boxrule=0.5pt, left=8pt, right=8pt, top=6pt, bottom=6pt,
  fonttitle=\bfseries\small, coltitle=white, colbacktitle=aimlsuccess,
  attach boxed title to top left={yshift=-2mm, xshift=4mm},
  boxed title style={boxrule=0pt, arc=1pt}, #1}

% ---- small semantics -------------------------------------------------
\newcommand{\scope}[1]{\textcolor{aimlmuted}{\small #1}}
\newcommand{\term}[1]{\textbf{#1}}
\newcommand{\code}[1]{\texttt{\small #1}}
\newcommand{\lecture}[1]{Lecture~#1}

\newtheorem{definition}{Definition}[section]
\newtheorem{proposition}[definition]{Proposition}

% ---- title block -----------------------------------------------------
% \notestitle{19}{Information retrieval: the lexical foundation}{SciFact (BEIR)}
\newcommand{\notestitle}[3]{%
  \thispagestyle{empty}
  \begin{flushleft}
    {\color{aimlmuted}\small\sffamily
     APPLICATIONS OF MACHINE LEARNING \quad$\cdot$\quad
     BSc Mathematics of Artificial Intelligence}\\[2mm]
    {\color{aimlrule}\rule{\textwidth}{0.8pt}}\\[4mm]
    {\color{aimlmuted}\large Lecture #1 \quad$\cdot$\quad Part V \quad$\cdot$\quad
     Extended notes}\\[3mm]
    {\Huge\bfseries\color{aimlprimary}\baselineskip=1.15\baselineskip #2\par}\vspace{4mm}
    {\color{aimlmuted}\large Dataset: #3}\\[3mm]
    {\color{aimlrule}\rule{\textwidth}{0.8pt}}
  \end{flushleft}
  \vspace{2mm}
  \begin{keybox}[title={Why these notes exist}]
    Lectures 19--22 sit outside G\'eron, the course's primary text. For those
    four lectures \emph{these notes are the primary source}, and they are
    examined on exactly the same terms as the textbook parts. Everything in
    the slide deck for this lecture is here, with the arguments written out at
    length rather than compressed onto a slide; the numbers are the ones the
    accompanying notebook prints, and every one of them is a measurement on
    the corpus named above rather than a figure quoted from a paper.
  \end{keybox}
  \vspace{1mm}
  {\small\tableofcontents}
  \vspace{2mm}
  \noindent{\color{aimlrule}\rule{\textwidth}{0.4pt}}
  \vspace{1mm}
}


\begin{document}
\notestitlefor{18}{IV}{Attention and transformers}{IMDb film reviews}{Chapters 14--15}

\section{One change removes three problems}

Lecture 16 ended with three things recurrence cannot do, and Lecture 17 built a
classifier that ran into all of them: the computation is sequential along the
sequence; everything remembered must fit through one fixed-width state; and
position 400 is four hundred steps from position 2, however good the gates.

Let every position compute a weighted sum over every position \emph{directly},
with the weights learned from the content. No recurrence, no single state, and
every pair of positions one step apart.

\section{Derivation: scaled dot-product attention}

The output at position $i$ is $\sum_j w_{ij}\mathbf{v}_j$ --- a weighted sum of
values, one per position. Note what is already gone: there is no recurrence in
that expression, so nothing forces the positions to be computed in order.
Everything below is about where $w_{ij}$ comes from.

Project each position three ways with learned matrices. A \term{query} is what
this position is looking for; a \term{key} is what that position offers; a
\term{value} is what it contributes if chosen. Score every pair by
$s_{ij} = \mathbf{q}_i^{\mathsf T}\mathbf{k}_j$ --- a dot product, because it
is the cheapest similarity a matrix multiply computes for all pairs at once ---
then normalise along the sequence with Lecture 17's softmax and take the
weighted sum. For all positions at once:
$\mathrm{softmax}(\mathbf{Q}\mathbf{K}^{\mathsf T})\mathbf{V}$.

\begin{derivbox}[title={Why divide by $\sqrt{d_k}$}]
Take $\mathbf{q}$ and $\mathbf{k}$ to have independent components with mean 0
and variance 1. Then $s = \sum_{d=1}^{d_k} q_d k_d$ is a sum of $d_k$
independent zero-mean terms, each of variance
$\operatorname{Var}(q_d)\operatorname{Var}(k_d) = 1$, so
\[ \operatorname{Var}(s) = d_k. \]
The scores grow with the dimension: at $d_k = 64$, typical scores are already
of order 8. Softmax of large scores saturates --- one weight goes to 1, the
rest to 0, and its gradient goes to zero with them, which is Lecture 11's
vanishing gradient arriving through a different door.

Dividing by $\sqrt{d_k}$ returns the variance to 1 whatever the dimension:
\[ \mathrm{Attention}
   = \mathrm{softmax}\!\left(
     \frac{\mathbf{Q}\mathbf{K}^{\mathsf T}}{\sqrt{d_k}}\right)\mathbf{V}. \]
\end{derivbox}

\emph{The scale is not a hyperparameter.} It is the number that makes the score
distribution independent of $d_k$.

\begin{keybox}[title={The trade, stated plainly}]
Every pair of positions interacts directly, so the path length between any two
is \textbf{1}, against $O(T)$ for a recurrent net --- and all positions compute
at once, so the sequence axis parallelises. The cost is $O(T^2)$ in time and
memory.

Recurrence is linear in $T$ and sequential, with a long gradient path.
Attention is quadratic in $T$ and parallel, with a path of length one. On the
sequence lengths people actually use the second wins --- and when it does not,
the $T^2$ term is what has to be approximated.
\end{keybox}

\subsection{Heads, and positions}

One attention layer computes one notion of relevance. Run $h$ of them on
different learned projections, each of width $d_{\text{model}}/h$ so that $h$
heads cost the same as one wide one, and concatenate. Different heads learn
different relations --- adjacency, agreement, coreference --- and nothing tells
them to; the division of labour is learned. That is Lecture 7's ensemble
argument in a new place: several estimators of the same thing beat one,
provided their errors are not identical, and different projections are what
stops them being identical.

\begin{failbox}[title={Failure condition --- the problem attention creates}]
$\mathrm{softmax}(\mathbf{Q}\mathbf{K}^{\mathsf T})\mathbf{V}$ is
\emph{permutation equivariant}: shuffle the input positions and the outputs
shuffle with them, unchanged.

So the model literally cannot tell ``the film was not good'' from ``good, the
film was not''. Recurrence got order for free by processing in order.
Attention has to be \emph{told}.
\end{failbox}

Positional encoding adds a position-dependent vector before the first layer:
\emph{sinusoidal}, fixed and extending to unseen lengths; \emph{learned}, an
embedding table indexed by position; or \emph{relative}, encoding the distance
between two positions in the score, which is what most current models do.

\subsection{One block, and what masking is for}

A transformer block is multi-head self-attention, then add the input back and
normalise, then a position-wise feed-forward network applied identically at
every position, then add and normalise again. Both wrappers are Lecture 11
unchanged: the residual gives the gradient an additive path through depth ---
the same argument as an LSTM gate, one level up --- and the normalisation keeps
the variance from drifting as blocks stack. \emph{A transformer is deep, and it
is deep only because of these two.}

A decoder generating text must not look at what it has not written, so
$s_{ij} = -\infty$ for $j > i$ before the softmax. $e^{-\infty} = 0$, so those
positions get exactly zero weight and contribute nothing --- and, crucially, no
gradient either. That is Lecture 15's causality condition again, enforced in
the architecture rather than in the data split: \emph{an unmasked decoder is a
model that has read the answer.}

\begin{center}
\small
\begin{tabular}{lll}
\toprule
\textbf{Family} & \textbf{Attention} & \textbf{Good at} \\
\midrule
Encoder-only & bidirectional & classification, tagging, retrieval \\
Decoder-only & masked, causal & generation \\
Encoder--decoder & both, plus cross-attention & translation, summarisation \\
\bottomrule
\end{tabular}
\end{center}

The choice is the same question as bidirectional layers in Lecture 16: at
prediction time, do I have the whole input? Today we classify a review we can
see in full, so encoder-only is right.

Self-attention and cross-attention are the \emph{same arithmetic} --- only
where $\mathbf{Q}$ and $\mathbf{K}$ come from changes, which is why one
implementation serves both.

\begin{keybox}[title={Where the parameters are}]
Attention is $4d_{\text{model}}^2$ per layer, for the four projections; the
feed-forward network is $8d_{\text{model}}^2$, being a fourfold expansion and
back.

\emph{Two thirds of the parameters in a transformer block are in the
feed-forward, not in the attention.} The attention is what makes it work and it
is the minority of the weights --- a useful corrective to the name.
\end{keybox}

Attention itself is six lines of code, and every one is a step of the
derivation. The rest of a transformer is arrangement around that function.
Which is worth saying plainly: \emph{the architecture that changed the field is,
at its centre, a similarity, a softmax and a weighted average.}

Implement it and check it against the library. The usual causes of
disagreement are a transposed key matrix, a mask applied after the softmax
instead of before, or a missing scale --- and each of those \emph{runs}. Only
the comparison finds them.

\section{Fine-tuning a pretrained transformer}

The model we borrow has 66{,}955{,}010 parameters, the same 30{,}522-token
vocabulary as last lecture, an embedding width of 768, and was pretrained to
predict masked words. \emph{Last lecture we took $30{,}522\times768$ numbers out
of this model and threw the rest away. The rest is where the language is.}

Fine-tuning puts a new, randomly initialised head on a pretrained body, at a
learning rate fifty times smaller than we used from scratch --- large steps
destroy what was pretrained --- for one pass over the data rather than four.

\begin{keybox}[title={How much is the head, and how much is the body?}]
The same backbone with its untrained head and no fine-tuning scores
\textbf{40.1\%}, against 50.0\% for always predicting one class, on 5{,}000 test
reviews of which 49.8\% are positive.

Note it is \emph{below} the majority class: an untrained head is an arbitrary
hyperplane and which side it puts each class on is a coin toss. The pretraining
on its own predicts nothing about sentiment, because nothing has told it what
the two output units mean.
\end{keybox}

\begin{center}
\small
\begin{tabular}{lr}
\toprule
\textbf{Model} & \textbf{Test accuracy} \\
\midrule
Always one class & 50.0\% \\
tf-idf + logistic regression & 90.1\% \\
GRU from scratch & 83.4\% \\
GRU, pretrained embeddings & 84.2\% \\
DistilBERT, fine-tuned & 90.6\% \\
\bottomrule
\end{tabular}
\end{center}

7.18 points over the model built last lecture --- and only 0.45 points over the
bag of words, which is the second result of the afternoon and the less
comfortable one.

\begin{keybox}[title={Say it in errors, not in points}]
The GRU misclassifies 4{,}146 reviews out of 25{,}000; the fine-tuned
transformer 2{,}352. \emph{43\% of the mistakes removed.}

Percentage points flatter a model near the top of the range. The error count is
what the desk cares about, because each one is a review someone has to fix by
hand.
\end{keybox}

\subsection{Where the pretrained model earns its place}

\begin{center}
\small
\begin{tabular}{rrrr}
\toprule
\textbf{Labelled reviews} & \textbf{Bag of words} & \textbf{DistilBERT} &
\textbf{Gap} \\
\midrule
   200 & 75.7\% & 83.9\% & 8.12 \\
 1{,}000 & 83.0\% & 87.0\% & 3.98 \\
 5{,}000 & 88.2\% & 88.9\% & 0.74 \\
20{,}000 & 90.5\% & 90.1\% & $-0.34$ \\
\bottomrule
\end{tabular}
\end{center}

The gap is largest where labels are scarcest, which is the situation every real
feedback desk starts in. \emph{The pretraining is worth thousands of labels},
and that is the number to quote to whoever is paying for the annotation.

Each model's budget is chosen on the validation split at every size --- the
number of optimisation steps for the transformer, the regularisation strength
for the linear model. Fixing the epochs instead would give the 200-review model
thirteen steps and fixing the steps would undertrain the 20{,}000-review one,
and measured, those two protocols disagree about the \emph{sign} at 20{,}000
labels.

\begin{center}
\small
\begin{tabular}{lrrr}
\toprule
\textbf{Model} & \textbf{Accuracy} & \textbf{Scoring 25{,}000} &
\textbf{Per review} \\
\midrule
tf-idf + logistic regression & 90.1\% & 3.6 s & 0.14 ms \\
GRU from scratch & 83.4\% & 6.3 s & 0.25 ms \\
DistilBERT, fine-tuned & 90.6\% & 133 s & 5.30 ms \\
\bottomrule
\end{tabular}
\end{center}

0.45 points for 37 times the cost per review --- \emph{a decision, not a
result}. At 20{,}000 labelled reviews the desk would probably make it against
us; at 200 it would not, and the table above says by how much. (The scoring
time is the least reproducible number here: five passes ran from 103 s to 160
s. Quote it to two figures or not at all.)

\section{The same representation, two more deliverables}

The desk wanted two more things --- ``find me the ones like this one'' and
``tell me what people keep complaining about''. Neither is classification, and
both fall out of the same pretrained models, because what those models produce
is a \emph{representation}, and a label was only ever one thing to do with it.

Mean-pooling over the \emph{real} tokens gives one vector per review --- and
that pooling is last lecture's padding bug in its other costume: divide by the
padded length instead of the true one and every short review is scaled towards
zero. Normalising onto the unit sphere makes the dot product the cosine, so
length carries no meaning and a long review cannot dominate every ranking.

Searching 5{,}000 negative reviews, embedded once in 11 seconds at 384
dimensions, is then one matrix product --- no index, no training, no labels.

\begin{center}
\small
\begin{tabular}{lr}
\toprule
\textbf{Query} & \textbf{Shared in the top 3} \\
\midrule
``the acting was wooden and unconvincing'' & 0 of 3 \\
``the sound mix made the dialogue impossible to follow'' & 0 of 3 \\
``far too long, it should have ended an hour earlier'' & 0 of 3 \\
\bottomrule
\end{tabular}
\end{center}

Not one document in common with keyword search, on any of the three. Keyword
search can only retrieve documents that reuse the query's words, and a
complaint rarely uses the desk's vocabulary.

\begin{failbox}[title={Failure condition --- that is a difference, not a victory}]
Nothing here measures which list is \emph{better}. Judged by eye, on one of
these three the keyword search returns the more relevant review. And cosine
similarity has no notion of negation, so \emph{the sound was perfect} and
\emph{the sound was not perfect} sit close together.

If you deploy this you owe yourself labelled queries and a retrieval metric ---
which is Lecture 19.
\end{failbox}

Clustering the same vectors is Lecture 8 unchanged, with $k$ chosen by
silhouette rather than by eye. Four groups, at a silhouette of 0.0179 --- low
in absolute terms, because complaint text does not fall into well-separated
balls, and the groups are still legible: 1{,}936 reviews around \emph{bad,
worst, movies, acting}; 1{,}198 around \emph{woman, character, women, role};
1{,}170 around \emph{horror, monster, zombie, scary}. Those terms are not the
cluster centres --- they are the words whose average tf-idf weight is highest
inside the group relative to outside it, because \emph{a group is named by what
makes it different, not by what is common everywhere}.

\section{One more prompt, and this is the one to remember}

\emph{``Build a scikit-learn pipeline that vectorises the reviews with tf-idf
and classifies them with logistic regression, and report the test accuracy.''}
Under-specified in the same place as Lecture 1's --- and this time the object
being fitted is not a scaler.

The code that comes back calls \code{fit\_transform} on \emph{every} document
and then splits. It runs, it splits, it never fits the classifier on test rows,
and it reports a plausible number.

\begin{failbox}[title={Failure condition --- reviewer question 2: what was fitted, and on what?}]
\code{TfidfVectorizer.fit} learns two things, and both come from the test rows
here: \emph{the vocabulary} --- which terms get a column at all, so terms
occurring only in test documents get their own feature --- and \emph{the
inverse document frequencies}, computed over the whole corpus.

A scaler learns two numbers per column from an invertible affine map. A
vectoriser learns \emph{which columns exist}, and that is not a transformation
of the features --- it is a choice of feature space, made with the answers in
the room.
\end{failbox}

\begin{center}
\small
\begin{tabular}{lrrrr}
\toprule
\textbf{Corpus} & \textbf{Leaky} & \textbf{Honest} & \textbf{Difference} &
\textbf{sd over seeds} \\
\midrule
25{,}000 reviews & 90.28\% & 90.22\% & 0.06 pts & 0.07\% \\
400 reviews & 75.90\% & 75.60\% & 0.30 pts & 2.53\% \\
\bottomrule
\end{tabular}
\end{center}

Read the last column before the one beside it. On the small corpus the mean gap
is 0.30 points and the spread across seeds is eight times larger, and the leak
wins on 10 of 20 seeds --- exactly half. \emph{This is a null result.}

Why so cheap? The leaky vocabulary really is bigger, 60{,}000 columns against
54{,}074. But a term occurring only in test documents has an all-zero column in
the training rows, so logistic regression gives it coefficient exactly zero ---
there is no gradient to move it. What actually leaks is the inverse document
frequencies, and an idf is an average over the whole corpus, so removing a
quarter of the documents barely moves it. The leak changes the \emph{scale} of
features the model already had, and a regularised linear model is nearly
indifferent to that.

\begin{keybox}[title={And the decision rule is not about the size}]
Fit the vectoriser inside the pipeline anyway. Not because this leak is
expensive --- measured, it is not --- but because \emph{the next fitted object
may not have that property}, and you will not re-derive the argument each time.

Three leaks, three small numbers: Lecture 2's scaler was worth nothing
measurable, Lecture 10's per-batch mean 0.382 points, today's vectoriser
nothing you can measure. A leaked score and an honest score can be identical,
so you cannot tell from the number which one you are holding. The procedure is
what separates them.
\end{keybox}

\subsection{The text leak that really is expensive}

The same document on both sides of the split. No pipeline protects you: every
object is fitted correctly, on training rows only, and the test row is still
one the model has memorised. On scraped or user-submitted text it is the normal
case --- a resubmit button is enough.

Checked on the official IMDb split: \textbf{123} test reviews are also present
in training, and 96 reviews are duplicated within training. The corpus authors
separated the \emph{films} between train and test, and said so; they did not
deduplicate the reviews. Three lines to run, and you now know rather than
assume. \emph{It is also exactly the check nobody ran, on the corpus everybody
uses.}

123 out of 25{,}000 is too few to measure a cost with, so build a corpus where
it is not: 1{,}500 reviews with three in ten submitted twice, 1{,}950 rows, and
under a random split 35\% of test rows have a copy of themselves in training.

\begin{center}
\small
\begin{tabular}{lrr}
\toprule
\textbf{Split} & \textbf{Reported accuracy} & \textbf{Spread over 10 seeds} \\
\midrule
random over rows & 90.0\% & 1.87\% \\
grouped by entry & 82.17\% & 2.88\% \\
\bottomrule
\end{tabular}
\end{center}

7.83 points, and the leak wins on 10 of 10 seeds. Nothing was fitted on the
test set in either case --- \emph{the difference is entirely which rows the
split happened to separate}. Compare with the vectoriser leak: the object that
was fitted wrongly cost far less than the rows that were split wrongly.

\begin{center}
\small
\begin{tabular}{llr}
\toprule
\textbf{Failure} & \textbf{Diagnosed in} & \textbf{Measured cost} \\
\midrule
Scaler fitted before the split & Lecture 2 & nothing measurable \\
Accuracy under class imbalance & Lecture 3 & a useless model above 90\% \\
Missing \code{zero\_grad()} & Lecture 10 & 76.7 points \\
Random split on a time series & Lecture 15 & an unreproducible score \\
Summarising at the padded last position & Lecture 17 & 18.24 points \\
Softmax before the loss & today & 1.67 points \\
Vectoriser fitted before the split & today & 0.30 points, not distinguishable
                                             from zero \\
Duplicate documents across the split & today & 7.83 points \\
\bottomrule
\end{tabular}
\end{center}

\section{Is 90.6\% good?}

Against the desk's own criterion: 2{,}352 reviews out of 25{,}000 are routed
wrongly and nobody sees them again. Whether that is acceptable depends on what
happens to a missed complaint, and we have not asked.

\begin{keybox}[title={What a defensible report says}]
Quote the error count, not the percentage. Say which reviews are misread ---
short ones, sarcastic ones, ones whose verdict is in the last clause. Offer an
abstain band around a probability of one half, and let the desk choose how wide
it is.

\emph{The model is not the deliverable. The routing decision is.}
\end{keybox}

What carried over unchanged: compute a trivial and a serious baseline before
committing (Lecture 1); split first, and fitted objects see training rows only
(Lecture 2); choose $k$ by silhouette rather than by eye (Lecture 8); own the
loop, and count a metric over the set (Lecture 10); and a pretrained
representation beats a bigger architecture (Lecture 13).

\section{The notebook}

\begin{trybox}[title={What to change}]
Shuffle the token positions of an input and run it through attention with and
without positional encoding. Without, the output is the same multiset; with, it
is not --- permutation equivariance, demonstrated rather than asserted.
\end{trybox}

\section{Exercises}

\begin{enumerate}[leftmargin=1.6em,itemsep=4pt]
  \item Derive the variance of a dot-product attention score under the standard
        assumptions, and state precisely what dividing by $\sqrt{d_k}$ fixes.
        \hfill \scope{[5]}
  \item Explain why self-attention is permutation equivariant, and what follows
        for a model that must distinguish ``not good'' from ``good \ldots
        not''. \hfill \scope{[4]}
  \item A decoder is trained without a causal mask and reports excellent
        validation loss. State what it has learned to do, and which earlier
        lecture's failure this is. \hfill \scope{[4]}
  \item A \code{TfidfVectorizer} is fitted before the split. Name the two
        quantities that leak, and explain why the measured cost is nevertheless
        near zero on a large corpus. \hfill \scope{[5]}
  \item A corpus contains duplicate documents and is split at random. Explain
        why no pipeline prevents the resulting leak, and give the split that
        does. \hfill \scope{[4]}
\end{enumerate}

\section*{Summary}
\addcontentsline{toc}{section}{Summary}

Attention is three projections, a scaled dot product, a softmax and a weighted
sum --- and the scale is not a hyperparameter: the score variance is $d_k$, so
dividing by $\sqrt{d_k}$ is what keeps the softmax off its saturated tails
whatever the width. Every pair of positions is one step apart and the sequence
axis parallelises, at a cost quadratic in length. Two thirds of a block's
parameters are in the feed-forward rather than the attention.

Fine-tuning a 67-million-parameter pretrained body reached 90.6\%: 7.18 points
over the recurrent model, 43\% of its errors removed --- and only 0.45 points
over tf-idf with logistic regression, at 37 times the cost per review. The
pretraining is worth thousands of labels, and the gap closes as labels arrive:
8.12 points at 200 reviews, and $-0.34$ at 20{,}000.

Two leaks, measured against each other. A vectoriser fitted before the split
learns the vocabulary and the idfs from the test rows, and costs 0.30 points on
a small corpus with a seed spread eight times larger --- a null result, and the
rule stays procedural anyway, because the next fitted object may not be so
forgiving. Duplicate documents across the split cost 7.83 points and win on 10
of 10 seeds, with nothing fitted wrongly at all. The official IMDb split
contains 123 of them.

\end{document}
